\documentclass[a4paper,11pt,fleqn]{article}%fleqn pour aligner les equations à gauche
\usepackage{../../Styles/Cours}

\newcommand{\chnum}{Chapitre 17}
\newcommand{\chtitle}{Caractérisation d'un dipôle ohmique}
\newcommand{\dtype}{TP}

\begin{document}

\maketitle

%\section*{Partie 2 : Caractéristique d'une résistance}
\begin{figure}[h!]
	\caption{Montage à réaliser}
	\centering
	\begin{circuitikz}[scale=1.2,transform shape]
		\draw (0,0)
			to [rmeter, t=GBF] (0,3)
			to [rmeter, t= A](3,3)
			to [R, l=$R_1$] (3,0)
			--(0,0);
		\draw (3,3)
			--(6,3)
			to [rmeter, t =V] (6,0)
			--(3,0);
		\end{circuitikz}
\end{figure}

On fait varier la valeur de tension du générateur de 0 à 6V par pas de 0,5V en notant à chaque fois les valeurs d'intensité et de tension que l'on reporte dans le tableau suivant :

\begin{tabular}{|>{\centering\arraybackslash}m{.3\textwidth}| >{\centering\arraybackslash}m{.3\textwidth}|>{\centering\arraybackslash}m{.3\textwidth}|}
	\hline
	Tension du générateur (V)& Intensité (A)& Tension $U_{R1}$ (V)\\
	\hline
	&& \rule[.8cm]{0cm}{0cm}\\
	\hline
	&& \rule[.8cm]{0cm}{0cm}\\
	\hline
	&& \rule[.8cm]{0cm}{0cm}\\
	\hline
	&& \rule[.8cm]{0cm}{0cm}\\
	\hline
	&& \rule[.8cm]{0cm}{0cm}\\
	\hline
	&& \rule[.8cm]{0cm}{0cm}\\
	\hline
	&& \rule[.8cm]{0cm}{0cm}\\
	\hline
	&& \rule[.8cm]{0cm}{0cm}\\
	\hline
	&& \rule[.8cm]{0cm}{0cm}\\
	\hline
	&& \rule[.8cm]{0cm}{0cm}\\
	\hline
	&& \rule[.8cm]{0cm}{0cm}\\
	\hline
	&& \rule[.8cm]{0cm}{0cm}\\
	\hline
	&& \rule[.8cm]{0cm}{0cm}\\
	\hline
\end{tabular}

\begin{figure}[h!]
	\begin{minipage}{.3\linewidth}
		\caption{Caractéristique d'un dipôle}
		La caractéristique d'un dipôle est la relation existante entre l'intensité $I$ du courant traversant le dipôle et la tension $U$ aux bornes de celui-ci. On représente la caractéristique  d'un dipôle en traçant graphiquement la courbe : $U = f(I)$.
	\end{minipage}\hspace{.04\linewidth}
	\begin{minipage}{.66\linewidth}
		\caption{Représenter et modéliser un graphique}
		\lstinputlisting[language = Python, numbers=left, stepnumber=1, basicstyle=\footnotesize]{Ressources/2GT_Act1.py}

		Pour représenter un nuage de points dans le langage de programmation python, on utilise le script ci-dessus dans lequel les valeurs ont été remplacées par des "...".\\
		\begin{itemize}
			\item Lignes 1-3 : importation des bibliothèques.
			\item Lignes 6-7 : np.array() crée des lignes de tableau, contenant ici les valeurs des abscisses et les ordonnées des points.
			\item Ligne 10 : plt.plot(x,y) trace le graphique de $y=f(x)$.
			\item Ligne 13 : linegress(x,y) calcule le coefficient directeur et l'ordonnée à l'origine de la droite de régression.
		\end{itemize}
	\end{minipage}
\end{figure}

\newpage
\subsection*{Questions :}
\begin{enumerate}
	\item Réaliser le montage et prendre les mesures indiquées dans le Document 1.
	\item Sur votre copie représenter graphiquement $I=f(U)$ à partir des points obtenus.
	\item Que peut on dire des points obtenus ?
	\item Ouvrir le fichier $"2GT\_Act1.py$" et compléter les "...".
	\item Relever les valeurs obtenues pour le coefficient directeur et l'ordonnée à l'origine de la droite de régression.
	\item Quelle loi permet de vérifier cette courbe caractéristique ?
\end{enumerate}
\end{document}

\section*{Partie 1 : Loi des nœuds, loi des mailles et loi d'Ohm}

\subsubsection*{Document 1 : Schéma électrique}

\begin{circuitikz}
\draw (0,0)
	to [rmeter, t=G, v>=$U_G$, i=$I$] (0,3) 
	to [R,l=$R_1$, i=$I_1$] (3,3)
	to [R,l=$R_2$ , i=$I_2$] (3,0) 
	-- (0,0);
\node [above] at (0,3) {A};
\node [above] at (0.7,3) {B};
\node [above] at (3,3) {C};
\node [above] at (6,3) {D};
\node [below] at (6,0) {E};
\node [below] at (3,0) {F};
%\draw [help lines] (-6,-6) grid (6,6);
\draw (3,3)
	to [R,l=$R_3$] (6,3)
	to [R,l=$R_4$, i=$I_3$] (6,0)
	-- (3,0);
\end{circuitikz}

\subsubsection*{Document 2 : Représenter une tension}
Comment flécher une tension électrique aux bornes d'un dipôle pour que la tension soit positive ?

\begin{circuitikz}
\draw (0,0)
	to [R, v=$U_{AB}$, i=$I$, l = R] (3,0);
\node [above] at (0,0) {A};
\node [above] at (3,0) {B};
\end{circuitikz}

\subsubsection*{Document 3 : Matériel à sélectionner}
\begin{itemize}
	\item Générateur de tension 6V
	\item Dipôles ohmiques de valeurs $R_1 = R_3 = 68 \Omega$ et $R_2 = R_4 = 100 \Omega$
	\item Un voltmètre et un ampèremètre 
\end{itemize}
\subsubsection*{Questions :}
\begin{enumerate}
		\item Flécher les tensions $U_1$, $U_2$, $U_3$ et $U_4$ aux bornes des résistances. Déterminer le nombre de mailles du circuit, les nommer et donner le nom des nœuds électriques.
		\item Réaliser le circuit du document 1. en respectant les valeurs des résistances indiquées.
		\item Rappeler le mode de branchement d'un voltmètre et d'un ampèremètre. Représenter un voltmètre mesurant $U_4$ et un ampèremètre mesurant $I_1$.
		\item Mesurer toutes les grandeurs (tensions et intensités) du montage à l'aide de l'appareil de mesure adapté. Rassembler toutes les mesures dans un tableau.
		\item Compléter la phrase suivante correspondant à la loi des nœuds : à un nœuds la somme des courants entrants est égale à ... .
		\item Compléter la phrase suivante correspondant à la loi des mailles : dans une maille la somme des tensions des dipôles est égale à ... .
		\item En utilisant la loi d'Ohm, calculer les valeurs des résistances du circuit et les comparer avec les valeurs marquées sur celles-ci.
\end{enumerate}
